\section{Proofs of theoretical results}
\label{app:theory-proofs}

The structural proofs use the fixed active positions and neuron
correspondence of Section~\ref{sec:theory-observations}.
All active code values are nonzero and neuron scales are positive.
The grouping under consideration uses every one of its code labels.

\subsection{Observable dimension and calibration scales}
\label{app:theory-dimension}

\begin{proof}[Proof of Proposition~\ref{prop:theory-dimension}]
    Set
    $\mathcal V_{\mathbf I}
    :=\operatorname{range}[\mathbf U_{\mathbf I},-\mathbf N]$.
    For vectors $\mathbf a\in\mathbb R^m$ and
    $\mathbf v\in\mathbb R^H$, the equation
    $\mathbf U_{\mathbf I}\mathbf a-\mathbf N\mathbf v=\mathbf0$
    requires $a_{j_{\mathbf I}(p)}=v_{h(p)}$ along every edge.
    These coordinates are constant on each component of
    $G_{\mathbf I}$, and each component allows one independent
    constant. Rank--nullity gives
    $\dim\mathcal V_{\mathbf I}=m+H-g_{\mathbf I}$.

    The centering matrix $\boldsymbol\Pi$ is the orthogonal projector
    onto the complement of $\operatorname{range}(\mathbf N)$.
    Since this range is contained in $\mathcal V_{\mathbf I}$,
    \begin{equation}
        \mathcal V_{\mathbf I}
        =\mathcal S_{\mathbf I}
            \mathbin{\oplus}\operatorname{range}(\mathbf N).
        \label{eq:appendix-theory-space-decomposition}
    \end{equation}
    The sum is orthogonal and $\mathbf N$ has rank $H$.
    Hence $\dim\mathcal S_{\mathbf I}=m-g_{\mathbf I}$.

    An exact calibration has
    $\sigma_{\mathbf I,p}c_{j_{\mathbf I}(p)}
    =s_{h(p)}r_p$.
    Let $\Delta\mathbf c$ and $\Delta\mathbf s$ denote perturbations
    of the code values and scales. Divide each homogeneous
    perturbation equation by the corresponding nonzero exact weight:
    \begin{equation*}
        \frac{\Delta c_{j_{\mathbf I}(p)}}{c_{j_{\mathbf I}(p)}}
        =\frac{\Delta s_{h(p)}}{s_{h(p)}}.
    \end{equation*}
    Thus the numerical calibration system also has exactly one
    common scale freedom per component. Fixing one scale in each
    component makes its exact solution unique.
    Exponentiating a common shift of the logarithmic variables
    gives the positive rescaling stated in the proposition.
\end{proof}

The rows of $[\mathbf U_{\mathbf I},-\mathbf N]$ are the transposes
of the oriented incidence columns of $G_{\mathbf I}$.
Its left nullspace is the cycle space, of dimension
$E-m-H+g_{\mathbf I}$. A fundamental cycle gives the alternating
logarithmic sum, or equivalently the product relation, described
in Section~\ref{sec:theory-observations}.
For a forest this dimension is zero, so the logarithmic incidence
map is onto $\mathbb R^E$. Every positive collection of edge
magnitudes can therefore be fitted.

\subsection{Pair signatures and recovery}
\label{app:theory-pairs}

\begin{lemma}[Information in an ordered-pair signature]
    \label{lem:theory-signatures}
    For the basis and signatures of
    \eqref{eq:theory-signature},
    \begin{equation}
        \begin{aligned}
            \boldsymbol\delta_{pq}=\mathbf0
            &\ \Longleftrightarrow\
                j_{\mathbf I}(p)=j_{\mathbf I}(q),\\
            \boldsymbol\delta_{pq}=\boldsymbol\delta_{uv}
                \ne\mathbf0
            &\ \Longleftrightarrow\
                (j_{\mathbf I}(p),j_{\mathbf I}(q))
                =(j_{\mathbf I}(u),j_{\mathbf I}(v)).
        \end{aligned}
        \label{eq:appendix-theory-signature-information}
    \end{equation}
    Each pair is taken within one neuron.
\end{lemma}

\begin{proof}
    There is a full-row-rank matrix
    $\mathbf T\in\mathbb R^{r\times m}$ such that
    $\mathbf M_{\mathbf I}=\mathbf Q\mathbf T$.
    Write $\mathbf e_j^{(m)}$ for the $j$th standard basis vector
    of $\mathbb R^m$. Centering subtracts the same row within a
    neuron, so
    \begin{equation*}
        \boldsymbol\delta_{pq}\mathbf T
        =(\mathbf e_{j_{\mathbf I}(p)}^{(m)})^{\mathsf T}
            -(\mathbf e_{j_{\mathbf I}(q)}^{(m)})^{\mathsf T}.
    \end{equation*}
    Multiplication of a row by $\mathbf T$ is injective.
    A difference of two standard basis vectors is zero exactly
    when their indices agree. Two nonzero differences agree
    exactly when both their positive and negative indices agree.
    These observations prove both statements.
\end{proof}

\begin{proof}[Proof of Theorem~\ref{thm:theory-pair-recovery}]
    By Lemma~\ref{lem:theory-signatures}, every merge made by the
    rule joins positions of the same true code.
    Within one neuron all occurrences of a code are merged.
    Between two neurons, direct merging of occurrences of code $j$
    occurs exactly when the neurons share another code.
    Consequently, the transitive closure joins exactly the
    occurrences in each connected component of $W_{\mathbf I}(j)$.
    The assumed connectivity recovers every entire group.

    A candidate with the same contrast space has the same
    signatures and must satisfy all these merges.
    Its partition is therefore a coarsening of the recovered
    partition. Since both have $m$ nonempty groups, the partitions
    agree up to their group names.
\end{proof}

\begin{proof}[Proof of Corollary~\ref{cor:theory-task-diversity}]
    A grouping explains all magnitudes exactly when the columns
    of $\mathbf Z$ lie in its contrast space.
    For sufficiency of this characterization, a fitted centered
    logarithm leaves a residual in $\ker\boldsymbol\Pi
    =\operatorname{range}(\mathbf N)$, which can be absorbed by
    the logarithms of positive neuron scales.
    Since $\mathbf Z$ has rank $r$, its columns span the true
    contrast space. Any candidate space of dimension $r$
    containing them equals that space.
    Apply Theorem~\ref{thm:theory-pair-recovery}.
\end{proof}

\subsection{The exact cycle criterion}
\label{app:theory-cycles}

\begin{proof}[Proof of Theorem~\ref{thm:theory-cycle-equivalence}]
    Fundamental cycles form a basis of the left nullspace of
    $[\mathbf U_{\mathbf I},-\mathbf N]$. Therefore
    \begin{equation}
        \mathbf R_{\mathbf I}\mathbf U_{\mathbf I}=\mathbf0,
        \qquad
        \mathbf R_{\mathbf I}\mathbf N=\mathbf0,
        \qquad
        \ker\mathbf R_{\mathbf I}=\mathcal V_{\mathbf I},
        \label{eq:appendix-theory-cycle-kernel}
    \end{equation}
    where $\mathcal V_{\mathbf I}$ is defined in
    Appendix~\ref{app:theory-dimension}.
    Parallel edges contribute cycles of length two.

    If the contrast spaces agree, their orthogonal decompositions
    in \eqref{eq:appendix-theory-space-decomposition} give
    $\mathcal V_{\mathbf J}=\mathcal V_{\mathbf I}$.
    Every column of $\mathbf U_{\mathbf J}$ then belongs to
    $\ker\mathbf R_{\mathbf I}$, and the dimensions in
    \eqref{eq:theory-cycle-equivalence} agree.

    Conversely, $\mathbf R_{\mathbf I}\mathbf U_{\mathbf J}
    =\mathbf0$ places every column of $\mathbf U_{\mathbf J}$
    in $\mathcal V_{\mathbf I}$.
    Together with the columns of $\mathbf N$, this gives
    $\mathcal V_{\mathbf J}\subseteq\mathcal V_{\mathbf I}$,
    and centering gives
    $\mathcal S_{\mathbf J}\subseteq\mathcal S_{\mathbf I}$.
    The dimension condition makes this inclusion an equality.
\end{proof}

The criterion also describes a complete finite enumeration.
For a nonempty subset $A\subseteq\mathcal E$, let
$\mathbf x_A\in\{0,1\}^E$ be its indicator vector.
Call $A$ admissible if $\mathbf R_{\mathbf I}\mathbf x_A=\mathbf0$.
Partition the positions into admissible groups and retain the
partitions for which the number of groups minus the number of
calibration components equals $\dim\mathcal S_{\mathbf I}$.
Each column of a membership matrix is exactly one such indicator.
Theorem~\ref{thm:theory-cycle-equivalence} thus proves that this
construction enumerates precisely the magnitude-equivalent
partitions, up to code names.

\subsection{Sign recovery}
\label{app:theory-sign-proof}

\begin{proof}[Proof of Proposition~\ref{prop:theory-signs}]
    Assume $T\ge1$ and normalize each observed profile by its first
    coordinate, defining $\overline{\mathbf b}_p:=b_{p,1}\mathbf b_p$.
    The fixed position sign cancels, and
    $\overline{\mathbf b}_p
    =\epsilon_{j_{\mathbf I}(p),1}
        \boldsymbol\epsilon_{j_{\mathbf I}(p)}$.
    Distinct code profiles up to sign give exactly $m$ distinct
    normalized profiles. Their equality therefore recovers the groups.

    If $p_0$ is a representative of the group containing $p$,
    then $b_{p,1}b_{p_0,1}
    =\sigma_{\mathbf I,p}\sigma_{\mathbf I,p_0}$.
    This recovers the relative signs. A representation with $m$
    codes needs a separate code for each of the $m$ normalized
    profiles. The only remaining choices are their names and
    one common sign per group.
\end{proof}

\subsection{Selection on independent tasks}
\label{app:theory-transfer-proof}

\begin{proof}[Proof of Theorem~\ref{thm:theory-transfer}]
    Condition on the information used to construct the bank and
    fix the adaptation rules.
    For each candidate $\mathbf I$, its query loss averaged
    within task $\tau\in\mathcal H$ is an independent observation
    in $[0,L]$ across tasks, with mean
    $\mathcal R_{n,b}^{\mathrm{adapt}}(\mathbf I)$.
    For any $a>0$, Hoeffding's inequality and a union bound give
    \begin{equation*}
        \Pr\left\{
            \max_{\mathbf I\in\mathcal B}
            \left|\widehat{\mathcal R}_{\mathcal H}(\mathbf I)
                -\mathcal R_{n,b}^{\mathrm{adapt}}(\mathbf I)\right|
                >a
            \right\}
        \le 2|\mathcal B|\exp(-2T_{\mathcal H}a^2/L^2).
    \end{equation*}
    Set $a=L\sqrt{\log(2|\mathcal B|/\delta)/(2T_{\mathcal H})}$.
    On the complementary event, replacing empirical risks by
    their true risks costs at most $a$ on each side of the
    empirical minimization inequality. This proves the bound.
\end{proof}

\section{Additional theoretical results}
\label{app:theory-additional}

\subsection{Stability under errors in source directions}
\label{app:theory-noise}

Let $r=\dim\mathcal S_{\mathbf I}>0$ and define the largest group
size by
$n_{\max}:=\max_j|\{p:j_{\mathbf I}(p)=j\}|$.
Let $\mathbf P_{\mathbf I}$ be the orthogonal projector onto
$\mathcal S_{\mathbf I}$.
Suppose an estimated space $\widehat{\mathcal S}$ also has rank $r$,
with projector $\widehat{\mathbf P}$.
Let $\eta\ge0$ bound their operator-norm difference.

\begin{proposition}[Stable pair recovery]
    \label{prop:theory-stable-pairs}
    Suppose
    \begin{equation}
        \|\widehat{\mathbf P}-\mathbf P_{\mathbf I}\|_{\mathrm{op}}
        \le\eta<\frac1{8\sqrt{n_{\max}}}.
        \label{eq:appendix-theory-space-error}
    \end{equation}
    Compute signatures from an orthonormal basis of
    $\widehat{\mathcal S}$.
    Using $2\sqrt2\eta$ as the threshold for a zero signature and
    $4\sqrt2\eta$ as the threshold for equal signatures recovers
    exactly the merges of the noiseless rule.
    Under the assumptions of
    Theorem~\ref{thm:theory-pair-recovery}, this recovers all groups.
\end{proposition}

\begin{proof}
    Write $\mathbf M_{\mathbf I}=\mathbf Q\mathbf T$ as in
    Lemma~\ref{lem:theory-signatures}.
    Since centering is an orthogonal projection,
    \begin{equation*}
        \|\mathbf T\|_{\mathrm{op}}
        =\|\mathbf M_{\mathbf I}\|_{\mathrm{op}}
        \le\|\mathbf U_{\mathbf I}\|_{\mathrm{op}}
        =\sqrt{n_{\max}}.
    \end{equation*}
    The image under $\mathbf T$ of a nonzero signature has norm
    $\sqrt2$. The difference of two distinct nonzero ordered-pair
    signatures has image of norm at least $\sqrt2$.
    Hence the distance from zero, and the separation between
    distinct nonzero signatures, are both at least
    $\sqrt{2/n_{\max}}$.

    The estimated orthonormal basis can be rotated to differ from
    $\mathbf Q$ in operator norm by at most
    $\varepsilon:=\sqrt2\eta$, by the principal-angle relation
    between equal-dimensional spaces.
    This rotation preserves signature distances.
    Each row has error at most $\varepsilon$, each signature
    at most $2\varepsilon$, and each difference of two signatures
    at most $4\varepsilon$.
    Condition~\eqref{eq:appendix-theory-space-error} makes the true
    separation strictly larger than $8\varepsilon$.
    The two stated thresholds therefore separate all required cases.
\end{proof}

This error criterion can be related to the observed directions.
Let the noiseless matrix $\mathbf Z$ of source contrasts have rank
$r$, let $\sigma_r(\mathbf Z)>0$ be its smallest nonzero singular
value, and write the observed matrix as
$\widetilde{\mathbf Z}=\mathbf Z+\mathbf E_Z$, where
$\mathbf E_Z\in\mathbb R^{E\times T}$ is its error.
Estimating the space by the first $r$ left singular vectors gives
\begin{equation}
    \|\widehat{\mathbf P}-\mathbf P_{\mathbf I}\|_{\mathrm{op}}
    \le\frac{2\|\mathbf E_Z\|_{\mathrm{op}}}
             {\sigma_r(\mathbf Z)}.
    \label{eq:appendix-theory-subspace-perturbation}
\end{equation}
Indeed, the residual of $\widetilde{\mathbf Z}$ outside its
estimated rank-$r$ space is at most $\|\mathbf E_Z\|_{\mathrm{op}}$.
The residual of $\mathbf Z$ is at most twice that amount.
Multiplying by its pseudoinverse proves the displayed bound.
For spaces of equal rank, the residual projector norm equals
the norm of the difference of projectors.

To control the logarithmic error, let $a>0$ satisfy
$|r_{\tau,p}|\ge a$ on all active positions and tasks.
Let $\Delta\mathbf R_{\mathrm{obs}}\in\mathbb R^{E\times T}$
collect their direction errors, each of magnitude at most $a/2$.
The logarithm is Lipschitz with constant $2/a$ on the relevant
magnitude intervals, so
\begin{equation}
    \|\mathbf E_Z\|_{\mathrm{op}}
    \le\|\mathbf E_Z\|_{\mathrm F}
    \le\frac2a\|\Delta\mathbf R_{\mathrm{obs}}\|_{\mathrm F}.
    \label{eq:appendix-theory-log-error}
\end{equation}
Together, \eqref{eq:appendix-theory-subspace-perturbation} and
Proposition~\ref{prop:theory-stable-pairs} give the sufficient
condition
\begin{equation*}
    \|\mathbf E_Z\|_{\mathrm{op}}
    <\frac{\sigma_r(\mathbf Z)}{16\sqrt{n_{\max}}}.
\end{equation*}
The denominator captures group size, while the singular value
captures excitation by the source tasks.
Small active weights amplify logarithmic errors.
Using this condition for fitted models requires an error bound
for their aligned directions.

\subsection{Uniqueness from longer cycles}
\label{app:theory-long-cycles}

Merge repeated occurrences of a code within each neuron and retain
one position from each resulting local group.
Lemma~\ref{lem:theory-signatures} determines these local groups from
the contrast space. Their representatives give a simple bipartite
calibration graph, denoted $G_{\mathbf I}^*$.
A graph is \emph{3-vertex-connected} if it has at least four vertices
and remains connected after deletion of any two or fewer vertices.

\begin{proposition}[A sufficient condition using all cycles]
    \label{prop:theory-three-connected}
    If $G_{\mathbf I}^*$ is 3-vertex-connected, every grouping
    $\mathbf J$ on the fixed active positions and neurons, with
    $m$ used codes and $\mathcal S_{\mathbf J}=\mathcal S_{\mathbf I}$,
    has the same groups up to code renaming.
    This follows from the cycle characterization and Whitney's
    graph-isomorphism theorem; see \citet{fomin2020whitney}.
\end{proposition}

\begin{proof}
    Equal contrast spaces force equal local groups by
    Lemma~\ref{lem:theory-signatures}. Restrict both models to
    their common representatives.
    Equation~\eqref{eq:appendix-theory-space-decomposition}
    gives equal incidence column spaces, hence equal cycle
    spaces on the same labelled edges.
    Minimal nonempty supports of vectors in a cycle space are
    the simple cycles, so the identity on edges preserves cycles.

    All edges of the 3-vertex-connected graph belong to one
    2-connected block. Cycle preservation places all candidate
    edges in one such block as well.
    Every candidate vertex is used, so its graph is 2-connected.
    Whitney's theorem then makes this edge bijection arise from
    a vertex isomorphism.
    For every known neuron, the incident edge set is fixed.
    Its at least two distinct edges have that neuron as their
    unique common endpoint in a simple graph.
    The isomorphism therefore fixes the neurons and only
    renames the code vertices.
\end{proof}

This condition can exploit cycles beyond the repeated pairs used
in Theorem~\ref{thm:theory-pair-recovery}.
Each is a sufficient condition in its own setting.
For graphs with small separators, candidate transformations must
also preserve the real oriented cycle space and the known neuron
blocks; preservation of cycle supports alone is weaker.

\subsection{Magnitude equivalence and different function classes}
\label{app:theory-counterexample}

Consider two unsigned assignments on two inputs and four neurons:
\begin{equation}
    \mathbf I=
    \begin{pmatrix}
        1&1&1&1\\
        2&2&3&3
    \end{pmatrix},
    \qquad
    \mathbf J=
    \begin{pmatrix}
        1&1&2&2\\
        3&3&3&3
    \end{pmatrix}.
    \label{eq:appendix-theory-ambiguous-pair}
\end{equation}
Both calibration graphs are connected and use three codes.
Identifying a vector of active positions with a $2\times4$ matrix,
their common contrast space consists of
\begin{equation*}
    \left\{
    \begin{pmatrix}
        u&u&v&v\\
        -u&-u&-v&-v
    \end{pmatrix}
    :u,v\in\mathbb R
    \right\}.
\end{equation*}
Thus known code count and graph connectivity alone allow ambiguity.

\begin{proposition}[Positive codes and signed separation]
    \label{prop:theory-functional-counterexample}
    Restricted to strictly positive code values, the assignments in
    \eqref{eq:appendix-theory-ambiguous-pair} define the same
    ReLU predictor class with free hidden biases and output layer.
    Allowing arbitrary real code values gives different classes.
\end{proposition}

\begin{proof}
    Equality of contrast spaces implies equality of the full
    logarithmic incidence spaces in
    \eqref{eq:appendix-theory-space-decomposition}.
    Any positive weights represented by one assignment can
    therefore be represented by the other up to positive neuron
    scales. Rescaling hidden biases and output weights
    accordingly preserves the predictor, in both directions.

    For separation with real codes, use
    $(c_1,c_2,c_3)=(1,1,-1)$ in $\mathbf I$.
    For $\mathbf x=(x_1,x_2)^{\mathsf T}$, define
    \begin{equation*}
        \begin{aligned}
            t_1(\mathbf x)&=x_1+x_2,&
            t_2(\mathbf x)&=x_1+x_2+1,\\
            t_3(\mathbf x)&=x_1-x_2+2,&
            t_4(\mathbf x)&=x_1-x_2+3.
        \end{aligned}
    \end{equation*}
    The function $f(\mathbf x):=\tfrac12\sum_{h=1}^4
    \sigma(t_h(\mathbf x))$ belongs to the class for $\mathbf I$.
    It has four distinct lines with nonzero gradient jumps.
    Any representation by four ReLU neurons must assign exactly
    one neuron to each line. After matching the lines, each
    contribution has the form $\tfrac12\sigma(t_h)$ or
    $\tfrac12\sigma(-t_h)$: its jump fixes the product of
    input scale and output coefficient.

    Let $F\subseteq\{1,\ldots,4\}$ contain the reversed
    contributions. Their change is
    $-\tfrac12\sum_{h\in F}t_h$, using
    $\sigma(-t)-\sigma(t)=-t$.
    For nonempty $F$, its coefficient of $x_1$ is
    $-|F|/2$, which an output constant cannot compensate.
    All four orientations must therefore agree with the
    displayed ones. Their second input weights have signs
    $+,+,-,-$. In $\mathbf J$ these weights share one code,
    so this orientation pattern is impossible.
    Hence $f$ belongs to the class for $\mathbf I$ and not
    to that for $\mathbf J$.
\end{proof}

The positive-code statement concerns strictly positive values;
its proof uses their logarithms.
The signed separation is a direct functional example with fixed
width and free biases and readout.

\subsection{Sign diversity and complete sign information}
\label{app:theory-sign-details}

Assume code signs are independent and balanced across codes and
tasks. Normalizing the first coordinate leaves $2^{T-1}$ equally
likely profiles. A union bound gives
\begin{equation}
    \Pr\{\text{a pair of code profiles agrees up to sign}\}
    \le \binom m2\,2^{1-T}.
    \label{eq:appendix-theory-sign-collisions}
\end{equation}
For $m\ge2$ and $\delta\in(0,1)$, a sufficient task count for this
probability to be at most $\delta$ is
\begin{equation*}
    T\ge1+\left\lceil
        \log_2\frac{\binom m2}{\delta}
        \right\rceil.
\end{equation*}
This statement isolates sign diversity; the observation model
still assumes correct support and neuron correspondence.

For noisy profiles, let $d_H$ denote Hamming distance and define
distance up to a common sign by
\begin{equation*}
    d_\pm(\mathbf u,\mathbf v)
    :=\min\{d_H(\mathbf u,\mathbf v),
                  T-d_H(\mathbf u,\mathbf v)\}.
\end{equation*}
Suppose at most $s$ signs in each position profile are wrong,
where $s$ is a nonnegative integer and $T>4s$.
If different code profiles are separated by more than $4s$
in $d_\pm$, connecting observed profiles at distance at most
$2s$ recovers exactly the true groups.
Indeed, within-group distances are at most $2s$, whereas
between-group distances exceed $2s$ by the triangle inequality.
The condition on $T$ also distinguishes the two possible
relative orientations within each group.

There is a corresponding exact description when every combination
of nonzero real code values is allowed.
Let $\mathbb F_2$ be the field with two elements, and let
$\mathbf a_{\mathbf I}\in\mathbb F_2^E$ record the fixed position
signs, with bit zero for $+1$ and bit one for $-1$.
The set of allowed sign vectors is the affine binary space
\begin{equation}
    \mathcal A_{\mathbf I}
    :=\mathbf a_{\mathbf I}
        +\operatorname{range}_{\mathbb F_2}(\mathbf U_{\mathbf I}).
    \label{eq:appendix-theory-sign-space}
\end{equation}

\begin{proposition}[Complete oriented weight information]
    \label{prop:theory-complete-signs}
    With fixed active positions and neurons, two models have the
    same families of nonzero weights modulo positive neuron
    scales exactly when both their contrast spaces and their
    affine sign spaces agree.
    Equality of the affine sign spaces determines the grouping
    and its relative signs up to code renaming and group sign changes.
\end{proposition}

\begin{proof}
    A nonzero weight vector is specified by its log magnitudes
    and signs. Positive neuron rescaling changes only the
    logarithmic part. Code magnitudes and signs vary independently,
    so the attainable pairs are exactly
    $\mathcal S_{\mathbf I}\times\mathcal A_{\mathbf I}$.

    Differences of elements of $\mathcal A_{\mathbf I}$ form
    $\operatorname{range}_{\mathbb F_2}(\mathbf U_{\mathbf I})$.
    Two positions have equal coordinates in every vector of
    this linear space exactly when they have the same code:
    different codes allow a sign change of only one of them.
    This recovers the partition.
    Within a group, the sum modulo two of two coordinates of
    any vector in $\mathcal A_{\mathbf I}$ is constant and
    gives their relative sign. The remaining choice is one
    sign per code.
\end{proof}

\subsection{Constructing normalization constraints}
\label{app:theory-normalization}

For one source task, let $r_0$ be the number of calibration variables
in the affine descriptor of Section~\ref{sec:method-projection}.
Its numerical variable is
$\boldsymbol\xi=(\mathbf c^{\mathsf T},
\mathbf a^{\mathsf T})^{\mathsf T}\in\mathbb R^{K+r_0}$.
Let $n_0:=K+r_0$ and let $\Omega\subseteq\mathbb R^{n_0}$ be an
open allowed region.
Let $q\ge1$ be the number of transformation parameters.
Suppose known smooth transformations on $\Omega$,
$\boldsymbol\gamma(\mathbf t,\boldsymbol\xi)$ preserve the
predictor and the current indexator, with
$\mathbf t\in\mathbb R^q$ and
$\boldsymbol\gamma(\mathbf0,\boldsymbol\xi)=\boldsymbol\xi$.
Compensating changes of unconstrained network parameters are
included in this assumption.
For a point $\boldsymbol\xi_0\in\Omega$, define the tangent matrix
\begin{equation*}
    \mathbf T_0
    :=\left.\frac{\partial\boldsymbol\gamma
        (\mathbf t,\boldsymbol\xi_0)}{\partial\mathbf t}
        \right|_{\mathbf t=\mathbf0}
    \in\mathbb R^{n_0\times q}.
\end{equation*}

\begin{proposition}[Local linear normalization]
    \label{prop:theory-local-normalization}
    Suppose $\mathbf T_0$ has rank $q$ and the transformations
    describe the local representation freedom.
    If $\mathbf G\in\mathbb R^{q\times n_0}$ satisfies
    $\det(\mathbf G\mathbf T_0)\ne0$, set
    $\mathbf h:=\mathbf G\boldsymbol\xi_0$.
    Every sufficiently close state has a unique sufficiently
    small transformation taking it to
    $\mathbf G\boldsymbol\xi=\mathbf h$.
    In particular, one can choose
    $\mathbf G=(\mathbf T_0^{\mathsf T}\mathbf T_0)^{-1}
        \mathbf T_0^{\mathsf T}$.
\end{proposition}

\begin{proof}
    Apply the implicit function theorem to
    $\mathbf G\boldsymbol\gamma(\mathbf t,\boldsymbol\xi)
    -\mathbf h=\mathbf0$ at $(\mathbf0,\boldsymbol\xi_0)$.
    Its derivative in $\mathbf t$ is the invertible matrix
    $\mathbf G\mathbf T_0$. The specified left inverse makes
    that matrix the identity.
\end{proof}

For diagonal positive scalings there is also a global construction
on a fixed sign region, related to the scaling sections studied by
\citet{hubert2012scalings}.
Let $\mathbf W\in\mathbb R^{n_0\times q}$ have rank $q$ and define
$\gamma_i(\mathbf t,\boldsymbol\xi)
:=\exp(\mathbf W_{i,:}\mathbf t)\xi_i$.
Assume $\Omega$ is preserved by these transformations.
Choose $q$ nonzero coordinates with ordered index set $S$
so that $\mathbf W_{S,:}$ is invertible.
Fix their signs throughout $\Omega$, and choose target values
$\mathbf h\in\mathbb R^q$ of the same signs.

\begin{proposition}[Global normalization for positive scalings]
    \label{prop:theory-global-normalization}
    Setting $\boldsymbol\xi_S=\mathbf h$ selects exactly one
    representative of each scaling orbit in $\Omega$.
    The required transformation is
    \begin{equation*}
        \mathbf t
        =\mathbf W_{S,:}^{-1}
            \bigl(\log|\mathbf h|-\log|\boldsymbol\xi_S|\bigr).
    \end{equation*}
\end{proposition}

\begin{proof}
    The signs are compatible. Taking log magnitudes of the
    $q$ required coordinate equalities gives a square linear
    system with the displayed invertible matrix.
    Its unique solution defines allowed positive scales.
\end{proof}

The one-scale-per-component rule of
Proposition~\ref{prop:theory-dimension} is a special case.
More generally, on a positive ReLU scale region one may fix a
positive weighted sum of neuron scales in each component.
Such a sum is positive and is multiplied by that component's
common scale, giving a unique normalizing factor.

Geometric normalization and uniqueness of a numerical least-squares
update have separate conditions.
At a fixed indexator, let $\mathbf E$ be its calibration residual
matrix, let $\mathbf v$ be its fixed affine offset, and let
$\mathbf D$ be a fixed positive diagonal discrepancy matrix,
as in Section~\ref{sec:method-projection}.
For feasible constraints $\mathbf G\boldsymbol\xi=\mathbf h$,
the minimizer of
$\|\mathbf D^{-1}(\mathbf E\boldsymbol\xi-\mathbf v)\|_2^2$
is unique exactly when
\begin{equation*}
    \ker\mathbf E\cap\ker\mathbf G=\{\mathbf0\}.
\end{equation*}
Indeed, the quadratic form is strictly positive on every nonzero
feasible direction precisely under this condition.
Preservation of a robust projection score additionally requires
that its residual and discrepancy scale transform by the same
positive factor coordinatewise. Their ratio is then invariant.

\subsection{Source selection and candidate coverage}
\label{app:theory-selection}

For the source procedure in
Section~\ref{sec:method-selection}, condition on the source tasks,
all proposal data, the final bank and all fitted candidate predictors.
Let $q_{\mathrm{str}}\ge1$ be the number of new independent
structure-selection examples per source task.
Let $R_{\mathrm{src}}(\mathbf I)$ be the mean true query loss
of these fixed predictors on the $T$ source tasks, and let
$\widehat{\mathbf I}_{\mathrm{src}}$ minimize their empirical
structure-selection loss.

\begin{proposition}[Conditional source-selection bound]
    \label{prop:theory-source-selection}
    For loss in $[0,L]$ and $\delta\in(0,1)$, with conditional
    probability at least $1-\delta$,
    \begin{equation}
        R_{\mathrm{src}}(\widehat{\mathbf I}_{\mathrm{src}})
        \le \min_{\mathbf I\in\mathcal B}R_{\mathrm{src}}(\mathbf I)
            +2L\sqrt{\frac{\log(2|\mathcal B|/\delta)}
                              {2Tq_{\mathrm{str}}}}.
        \label{eq:appendix-theory-source-selection}
    \end{equation}
\end{proposition}

\begin{proof}
    Conditional on the fixed tasks and predictors, the
    $Tq_{\mathrm{str}}$ example losses are independent and
    bounded. Apply Hoeffding's inequality to each candidate,
    take a union bound, and use the empirical minimization
    argument in the proof of Theorem~\ref{thm:theory-transfer}.
\end{proof}

This risk averages predictions on the particular source tasks.
The independent unit in Theorem~\ref{thm:theory-transfer} is instead
a newly sampled task, with fresh numerical fitting.
The two bounds therefore address different sources of uncertainty.

For the latter setting, define the bank's adaptation approximation
error relative to the admissible indexators by
\begin{equation*}
    \varepsilon_{\mathrm{bank}}
    :=\min_{\mathbf I\in\mathcal B}
        \mathcal R_{n,b}^{\mathrm{adapt}}(\mathbf I)
        -\inf_{\mathbf I\in\mathcal I_K}
        \mathcal R_{n,b}^{\mathrm{adapt}}(\mathbf I).
\end{equation*}
Subtracting the infimum from
\eqref{eq:theory-transfer-bound} gives
\begin{equation*}
    \mathcal R_{n,b}^{\mathrm{adapt}}
        (\widehat{\mathbf I}_{\mathcal H})
        -\inf_{\mathbf I\in\mathcal I_K}
            \mathcal R_{n,b}^{\mathrm{adapt}}(\mathbf I)
    \le \varepsilon_{\mathrm{bank}}
        +2L\sqrt{\frac{\log(2|\mathcal B|/\delta)}
                         {2T_{\mathcal H}}}.
\end{equation*}
The algebraic constraints in
Theorem~\ref{thm:theory-cycle-equivalence} preserve every compatible
partition if all admissible continuations are retained.
A fixed-size bank needs an additional coverage argument.
For example, if $M$ candidates are indistinguishable by the information
used for pruning and a bank retains at most $B_{\max}<M$ of them,
then under a uniform conditional distribution of the true candidate,
coverage probability is at most $B_{\max}/M$.
This follows by counting the retained candidates, and also holds
after averaging over randomized pruning.

\subsection{Prediction error and numerical conditioning}
\label{app:theory-adaptation}

A linear model gives an explicit separation of approximation error
and estimation variance, as in the analysis of learned sharing in
\citet{yeh2022equivariance}.
Let $\mathbf w^\star\in\mathbb R^P$ be the true weights,
$\mathbf X\in\mathbb R^{n\times P}$ a fixed design satisfying
$\mathbf X^{\mathsf T}\mathbf X=n\mathbf{Id}_P$, and
$\mathbf y=\mathbf X\mathbf w^\star+\boldsymbol\varepsilon$,
where $\sigma>0$ is the noise standard deviation and
$\boldsymbol\varepsilon\sim
\mathcal N(\mathbf0,\sigma^2\mathbf{Id}_n)$.
Let an independent query input have second moment
$\mathbb E[\mathbf x\mathbf x^{\mathsf T}]=\mathbf{Id}_P$.
Before observing target data, fix an allowed weight space with
orthonormal basis $\mathbf U\in\mathbb R^{P\times r}$ and
projector $\mathbf P_U:=\mathbf U\mathbf U^{\mathsf T}$.
Let $\widehat{\mathbf w}_U$ be the least-squares estimator in that space.

\begin{proposition}[Bias and variance of a fixed constraint]
    \label{prop:theory-bias-variance}
    Under these assumptions,
    \begin{equation}
        \mathbb E\bigl[
            (\mathbf x^{\mathsf T}
                (\widehat{\mathbf w}_U-\mathbf w^\star))^2
            \bigr]
        =\|(\mathbf{Id}_P-\mathbf P_U)\mathbf w^\star\|_2^2
            +\frac{\sigma^2r}{n}.
        \label{eq:appendix-theory-bias-variance}
    \end{equation}
\end{proposition}

\begin{proof}
    Orthogonality of the design gives
    $\widehat{\mathbf w}_U=\mathbf P_U\mathbf w^\star
    +n^{-1}\mathbf P_U\mathbf X^{\mathsf T}\boldsymbol\varepsilon$.
    The deterministic error is orthogonal to the random part,
    whose covariance is $\sigma^2\mathbf P_U/n$.
    The query second moment turns prediction error into squared
    weight error. Taking its expectation proves the formula.
\end{proof}

Dense least squares has variance $\sigma^2P/n$ and zero
approximation error in this model. Thus a constraint improves
prediction when its squared approximation error is smaller than
$\sigma^2(P-r)/n$.
This exact calculation uses $n\ge P$, as required by the design
assumption. Nonlinear fitting with a free output layer has a
different geometry.

For numerical conditioning, consider a positive definite quadratic
loss with Hessian $\mathbf H\in\mathbb R^{P\times P}$.
Restricting it to an orthonormal basis
$\mathbf U\in\mathbb R^{P\times r}$, with $r\ge1$, gives
the Hessian $\mathbf H_U:=\mathbf U^{\mathsf T}\mathbf H\mathbf U$.
The Rayleigh quotient yields
\begin{equation}
    \lambda_{\min}(\mathbf H_U)\ge\lambda_{\min}(\mathbf H),
    \qquad
    \lambda_{\max}(\mathbf H_U)\le\lambda_{\max}(\mathbf H).
    \label{eq:appendix-theory-conditioning}
\end{equation}
Its spectral condition number is therefore no larger.
These inequalities follow by minimizing and maximizing the
same quotient over a subset of the unit vectors.

The coordinates of an ordinary codebook need not be orthonormal.
Let $\mathbf L_{\mathbf I}\in\mathbb R^{P\times m}$ be the
signed lookup matrix, with active row
$\sigma_{\mathbf I,p}(\mathbf e_{j_{\mathbf I}(p)}^{(m)})^{\mathsf T}$
and zero row for an inactive position.
If $n_j$ is the size of group $j$, then
$\mathbf L_{\mathbf I}^{\mathsf T}\mathbf L_{\mathbf I}
=\operatorname{diag}(n_1,\ldots,n_m)$.
Even for $\mathbf H=\mathbf{Id}_P$, group sizes $1$ and an integer $N\ge1$
produce codebook condition number $N$.
The normalized lookup
$\mathbf L_{\mathbf I}\operatorname{diag}(n_1,\ldots,n_m)^{-1/2}$
has orthonormal columns and restores
\eqref{eq:appendix-theory-conditioning}.
This separates the geometry of a constraint from the effect of
its numerical coordinates. Training-time comparisons also depend
on initialization, nonlinear trajectories and the cost of each update.
